\documentclass[a4paper, 11pt, oneside]{report} 
\usepackage[utf8]{inputenc}
\usepackage[dutch]{babel}
\usepackage{amsmath}
\usepackage{amsfonts}
\usepackage{amssymb}
\usepackage{graphicx}
\usepackage{caption}
\usepackage[table,xcdraw]{xcolor}
\usepackage[toc,page]{appendix}
\usepackage{hyperref}
\usepackage{titlesec}
\usepackage{listings}
\usepackage{float}
\usepackage{tikz}
\usetikzlibrary{trees}
\usepackage{tikz-qtree}
\usepackage{graphicx}
\usepackage{fancyref}
\usepackage{wrapfig}
\usepackage{url}
\usepackage{pdflscape}
\usepackage{fancyvrb}
\usepackage{fancyhdr}
\graphicspath{ {Afbeeldingen/} }
\usepackage{subfig}
\usepackage{tabularx}
\usepackage{apacite}
\usepackage{longtable}
\usepackage{titlecaps}
%\usepackage[T1]{fontenc}
\usepackage{titlesec, blindtext, color}
\definecolor{gray75}{gray}{0.75}
\newcommand{\hsp}{\hspace{20pt}}
\usepackage{pdfpages}


\newcolumntype{L}[1]{>{\raggedright\arraybackslash}p{#1}}

\titleformat{\chapter}[hang]{\huge\bfseries}{\thechapter\hsp\textcolor{gray75}{|}\hsp}{0pt}{\Large\bfseries}


\def\sectionautorefname{Paragraaf}
\def\chapterautorefname{Hoofdstuk}
\def\tableautorefname{Tabel}
\DeclareRobustCommand{\VAN}[3]{#2} % set up for citation

%% Sets page size and margins 
\usepackage[a4paper,top=3cm,bottom=3cm,left=3cm,right=3cm,marginparwidth=1.75cm]{geometry}

\author{M.W.J. Berentsen}

\title{Onderzoeksrapport}
\usepackage{titling}

\newcommand{\subtitle}[8]{%
	\posttitle{%
		\par\end{center}
	\begin{center}\large#1\end{center}
	\vskip0.5em
	\begin{center}\large#2\end{center}
	\begin{center}\large#3\end{center}
	\begin{center}\large#4\end{center}
    \begin{center}\large#5\end{center}
    \begin{center}\large#6\end{center}
    \begin{center}\large#7\end{center}
    \begin{center}\large#8\end{center}
	\vskip0.5em}%
}

\subtitle{Drone mesh netwerk simulatie}{HAN Arnhem}{561399}{MWJ.Berentsen@student.han.nl}{Versie 1}{Alten Nederland B.V.}{Docent: J. Visch, MSc}{Assessor: ir. C.G.R. van Uffelen}

\setlength{\parindent}{0pt}
\setlength{\parskip}{5pt plus 2pt minus 1pt}



\hypersetup{colorlinks=true, urlcolor=red,citecolor=black,linkcolor=blue}  % Colours hyperlinks in blue, but this can be distracting if there are many links.
\setcounter{tocdepth}{2}



\begin{document}
\begin{figure}
\begin{center}\includegraphics[scale=0.1]{alten}\end{center}
\end{figure}
\maketitle

%\section*{Voorwoord}
%\addcontentsline{toc}{section}{\protect\numberline{}Voorwoord}
%\pagebreak

%Geschikt voor minimaal 50 nodes; Kan slecht of geen signaal nabootsen

\tableofcontents
\clearpage
%\section*{Begrippenlijst}

% Please add the following required packages to your document preamble:
% \usepackage[table,xcdraw]{xcolor}
% If you use beamer only pass "xcolor=table" option, i.e. \documentclass[xcolor=table]{beamer}
%\begin{table}[H]
%\centering

%\label{begrippen}
%\begin{tabular}{|l|l|}
%\hline
%\rowcolor[HTML]{C0C0C0}
%Term        & Omschrijving                                                         \\ \hline
%term        & Omschrijving                                                      	\\ \hline

%\end{tabular}
%\caption{Begrippenlijst}
%\end{table}

%\clearpage

%\section*{Samenvatting}
%\addcontentsline{toc}{section}{\protect\numberline{}Samenvatting}
%\pagebreak

\chapter*{Samenvatting}

\hrule
Optioneel een samenvatting van het onderzoek. Hier kunnen anderen snel inzicht krijgen in wat jij hebt onderzocht en wat je conclusie is.
\begin{itemize}
	\item Kunnen derden snel inzicht krijgen in jouw onderzoek?
	\item Staat de conclusie erin vermeld?
\end{itemize}

\chapter{Inleiding}
\label{chapter:inleiding}

\hrule
De inleiding beschrijft:
\begin{itemize}
\item Waarvoor het onderzoek gedaan wordt;
\item Beschrijf waarom het onderzoek nu wordt uitgevoerd;
\item Het doel van het onderzoek.
\end{itemize}



\chapter{Hoofd- en deelvragen}
In dit hoofdstuk worden de hoofd- en deelvragen genoemd en onderbouwd.
Er wordt een scope bepaalt met wat er onderzocht wordt en dus ook wat niet.
Vervolgens wordt de onderzoeksmethode toegelicht en beargumenteerd.
Tenslotte wordt de invloed van het onderzoek op het afstudeerproject beschreven.

\section{Hoofdvraag}
Het doel van dit onderzoeksrapport is het beantwoorden van de volgende opgesteld hoofdvraag:
\begin{quotation}
\textit{Is het mogelijk om meerdere op de grond geparkeerde drones, verbonden in een onderling mesh netwerk, in zetten voor het monitoren van gebieden?}	
\end{quotation}
Doordat de uitvoerende student van dit onderzoek niet beschikt over een certificaat om met drones te mogen vliegen behoudt dit onderzoek zich tot een simulatie.
Het mesh netwerk wordt wel voorzien van een fysiek prototype.

\section{Deelvragen}

Om de hoofdvraag te kunnen beantwoorden zijn er de volgende deelvragen opgesteld:

\begin{itemize}
	\item Welke simulatie software is geschikt om een drone valide en realistisch in te na te bootsen?
	\item Welke gesimuleerde drone is geschikt voor de simulatie software?
	\item Welke mesh netwerk hardware is geschikt voor het onderling verbinden van drones?
	\item Welke mesh netwerk libary is geschikt voor de te gebruiken hardware?
	\item Hoe gedraagt een node uit het mesh netwerk zich bij slecht of geen bereik?
\end{itemize}



\chapter{Criteria}

\hrule
Maak hier een lijst met criteria aan de hand van de hoofd- en deelvragen.
\begin{itemize}
	
\item Zijn de criteria duidelijk opgesteld?
\item Waar moet de oplossing aan voldoen, staat dit erin?
\end{itemize}

\begin{table}[H]
	\centering
	\begin{tabular}{|l|l|l|}
		\hline
		\rowcolor[HTML]{C0C0C0} 
		Naam & Beschrijving & Kosten \\ \hline
		&              &        \\ \hline
		&              &        \\ \hline
	\end{tabular}
	\caption{Opgestelde criteria}
	\label{tab:criteria}
\end{table}

\chapter{literatuur}

\hrule
Zet hier de achterliggende theorie.
\begin{itemize}
\item Staat erin vermeld dat je de bieb hebt toegepast uit de methodekaart van de HAN?
\item Staan er oplossingsrichtingen vermeld die voldoen aan de criteria?
\end{itemize}


\chapter{Oplossingsrichtingen}

\hrule
\begin{itemize}
\item Beschrijf je hoe jouw oplossingsrichting de hoofdvraag beantwoord?
\item Beschrijft de oplossingsrichting de mogelijkheden?
\item Staat hier "Veld" vermeldt van de methodekaart van de HAN?
\end{itemize}

\chapter{Experimenten}

\hrule
\begin{itemize}
\item Beschrijft waarom dit experiment van belang is.
\item Voldoet de code van het experiment aan de opgestelde eisen?
\item Voldoet het programma aan de opgestelde QoS?
\item Is gedocumenteerd waarom het programma voldoet aan de opgestelde eisen?
\item Is gedocumenteerd waarom het programma voldoet aan de opgestelde QoS?
\end{itemize}

\chapter{Conclusie}

\hrule
\label{chapter:conclusie}
Conclusie uit resultaten. Herhaal en beantwoord de hoofdvraag.
\begin{itemize}
\item Trek een conclusie uit de resultaten van de deelvragen.
\item Wordt de hoofdvraag beantwoord?
\item Wat beveel je aan?
\end{itemize}


\bibliographystyle{apacite}
\bibliography{bilbliography.bib}

\clearpage
\appendix

\end{document}